% !TeX program = lualatex
%
% dev 分支 / dev branch: 现代化重设计, 吸收 Moloch 与 Jacquenetta 的
% 设计语言, 与 main 分支的经典版主题并行维护.
% Modernized redesign on dev: design language absorbed from Moloch and
% Jacquenetta, maintained in parallel with the classic theme on main.
\documentclass[aspectratio=169]{beamer}

% 如果在 Overleaf 上编译，注意切换到 LuaLaTeX 作为编译器
% If you are using Overleaf, switch to LuaLaTeX as the compiler
%
% 推荐编译方式 / Recommended toolchain:
%   latexmk        （读取 .latexmkrc，自动执行 lualatex -> biber -> lualatex 直至收敛）
% 手动等价命令 / Manual equivalent:
%   lualatex slide.tex && biber slide && lualatex slide.tex && lualatex slide.tex
% 不推荐使用 bibtex，本项目按 biber 配置。
% bibtex is not recommended: this project is configured for biber.

% 字体设置 / Font settings
\usepackage{fontspec} % 字体选择 / Font selection
\usepackage{luatexja-fontspec} % 中文字体支持, LuaLaTeX 版 / Chinese font support for LuaLaTeX

% 设置中文字体 / Set Chinese fonts (luatexja-fontspec 宏名)
\setmainjfont{Noto Sans CJK SC}
\setsansjfont{Noto Sans CJK SC}
\setmonojfont{Noto Sans Mono CJK SC}

% 设置英文字体 / Set English fonts
\setsansfont{Fira Sans} % 正文默认无衬线 / Beamer body text is sans by default
\setmainfont{Fira Sans} % 主字体 / Main font
\setmonofont{Fira Mono} % 等宽字体 / Monospace font

% Overleaf 方案（无须自定义字体）/ Overleaf solution (no need to customize fonts)
% \setsansjfont{FandolHei}[
%   BoldFont = FandolHei-Bold
% ]
% \setmainjfont{FandolSong}[
%   BoldFont = FandolSong-Bold,
%   ItalicFont = FandolKai,
%   BoldItalicFont = FandolKai
% ]
% \setmonojfont{FandolFang}

% 参考文献 / References
\usepackage[style=numeric-comp, backend=biber]{biblatex} % 参考文献管理 / Bibliography management
\addbibresource{ref.bib} % 添加参考文献文件 / Add bibliography file

% 文档信息 / Document information
\author{篠原 はるな}
\title{SDU Beamer Theme}
\subtitle{山大风格 Beamer 模板与简易指南}
\institute{山东大学 软件学院}
\date{\today}

\usetheme{SDU} % 等价于旧写法 \usepackage{SDU} / equivalent to the legacy \usepackage{SDU}

% 让引用、目录、链接以主题色显示，便于辨认
% Color citations, TOC entries, and links in the theme color
\hypersetup{colorlinks=true, linkcolor=., citecolor=peaking, urlcolor=peaking!70!black}

\begin{document}

% 标题页 / Title page
\begin{frame}[plain]
    \titlepage
    \begin{figure}[htpb]
        \begin{center}
            \includegraphics[width=0.1\linewidth]{img/Example/SDU.eps}
        \end{center}
    \end{figure}
\end{frame}

% 目录页 / Table of contents
\begin{frame}[plain]
    \tableofcontents[sectionstyle=show,subsectionstyle=show/shaded/hide,subsubsectionstyle=show/shaded/hide]
\end{frame}

\section{模板简介}

\subsection{主题特性}

\begin{frame}{模板简介}
    本模板是面向山东大学的 \LaTeX{} Beamer 主题，适用于学术汇报、
    开题与答辩等场合，主要特性包括：
    \begin{itemize}
        \item 现代重设计：吸收 Moloch 与 Jacquenetta 的设计语言，
              与 \texttt{main} 分支的经典版并行维护
        \item 默认 \(16:9\) 宽屏长宽比，白底，山大红（CMYK 26/100/100/28）作强调色
        \item 无导航条，页脚为细进度条 + 页码；\texttt{[border]} 选项可加页面黑框
        \item 使用 Lua\LaTeX{} 编译，通过 \texttt{luatexja} 支持中文排版
        \item 配色集中在 \texttt{beamercolorthemeSDU.sty}，
              \texttt{\textbackslash sdusetaccent} 一条命令换强调色
        \item 每个 section 前自动生成深色章节页（幽灵数字 + 章节名）
        \item 三类 block 统一为左竖条 callout 风格，竖条颜色区分语义
        \item 预置代码高亮、三线表、参考文献等常用配置
    \end{itemize}
\end{frame}

\subsection{历史沿革}

\begin{frame}{两个版本，两条路径}
    本主题沿两条路径维护，谱系可追源到早期的清华 Beamer 主题%
    \footnote{最初版本发布在 latexstudio.net，原始下载链接现已失效。}：
    \begin{columns}[T]
        \begin{column}{0.5\textwidth}
            {\bfseries\color{peaking} \texttt{main}：经典版}
            \begin{itemize}
                \item THU-Beamer-Theme（TUNA 协会）\cite{thu}
                \item PKU-Beamer-Theme \cite{pku}
                \item SDU-Beamer-Theme（Ura 改配色与布局）\cite{origin}
                \item 本仓库重构：\(16:9\) 长宽比、配色抽离为
                      \texttt{beamercolorthemeSDU.sty}、GitHub Actions 自动编译
            \end{itemize}
        \end{column}
        \begin{column}{0.5\textwidth}
            {\bfseries\color{peaking} \texttt{dev}：现代版（本页所示）}
            \begin{itemize}
                \item 与经典版并行维护的现代化重设计
                \item 设计参考 Moloch\footnote{\url{https://github.com/jolars/moloch}}
                      与 Jacquenetta\footnote{\url{https://github.com/smanchajm/Jacquenetta-beamer-theme}}
                \item Lua\LaTeX{} + \texttt{luatexja} 引擎迁移
                \item 详见下一页
            \end{itemize}
        \end{column}
    \end{columns}
    \vspace{1em}
    校徽矢量图来自 \cite{sdulogo}。
\end{frame}

\begin{frame}{本分支：现代化重设计}
    你正在查看的这份文档由 \texttt{dev} 分支构建，是一次现代化重设计，
    与 \texttt{main} 分支的经典版主题并行维护：
    \begin{itemize}
        \item 设计语言：吸收 \textbf{Moloch}（Metropolis 的现代继承者）与
              \textbf{Jacquenetta} 的设计语言——白底、无导航条、细进度条、
              大字号无衬线标题、统一左竖条 callout、深色章节页
        \item 机构标识：仅保留山大红与校徽
        \item 引擎：从 Xe\LaTeX{} 迁移至 \textbf{Lua\LaTeX{}}，
              中文排版由 \texttt{luatexja} 处理
        \item 字体：Fira Sans 正文 + Noto Sans CJK 中文，
              中英文统一无衬线
        \item 需要经典版式请切换到 \texttt{main} 分支
    \end{itemize}
\end{frame}

\subsection{项目结构}

\begin{frame}{项目结构}
    \begin{table}[h]
        \centering
        \begin{tabular}{ll}
            \toprule
            文件                               & 作用                    \\
            \midrule
            \texttt{slide.tex}               & 主文档（本文件）              \\
            \texttt{beamerthemeSDU.sty}      & 主题样式：进度条、章节页、block 样式 \\
            \texttt{beamercolorthemeSDU.sty} & 配色定义，可整体替换            \\
            \texttt{ref.bib}                 & 参考文献条目                \\
            \texttt{img/}                    & 插图目录                  \\
            \texttt{fonts/}                  & 随仓库分发的中文字体            \\
            \bottomrule
        \end{tabular}
    \end{table}
\end{frame}

\subsection{编译方式}

\begin{frame}[fragile]{编译方式}
    需要支持中文的 Lua\LaTeX{} 工具链，本地编译命令为：
    \begin{lstlisting}[language=bash,numbers=none]
latexmk          # 读取 .latexmkrc，自动完成全部编译步骤
\end{lstlisting}
    手动等价为 \texttt{lualatex -> biber -> lualatex -> lualatex}：
    \begin{itemize}
        \item 使用 VS Code 可安装 \texttt{LaTeX Workshop} 插件，
              配置 \texttt{lualatex -> biber -> lualatex} recipe
        \item 使用 Overleaf 请将编译器切换为 LuaLaTeX
        \item 仓库自带 GitHub Actions，每次 push 自动编译并回传 PDF
    \end{itemize}
\end{frame}

\section{文本与列表}

\subsection{列表环境}

\begin{frame}{列表环境}
    \begin{columns}[t]
        \begin{column}{0.48\textwidth}
            \begin{itemize}[<+-| alert@+>]
                \item 无序列表：\texttt{itemize}
                \item 逐项出现并标红：\texttt{[<+-| alert@+>]}
                      \begin{itemize}
                          \item 支持二级嵌套
                      \end{itemize}
                \item 也可用 \texttt{\textbackslash item<2->} 手动控制
            \end{itemize}
            \begin{enumerate}
                \item 编号列表：\texttt{enumerate}
                \item 自动生成序号
            \end{enumerate}
        \end{column}
        \begin{column}{0.48\textwidth}
            \begin{description}
                \item[LaTeX] 专业排版系统
                \item[Beamer] 学术幻灯片文档类
                \item[LuaLaTeX] 原生支持 Unicode 字体的引擎
            \end{description}
        \end{column}
    \end{columns}
    \vspace{1em}
    \pause
    本行文字在第二步出现，由 \texttt{\textbackslash pause} 控制。
\end{frame}

\subsection{块环境}

\begin{frame}{块环境}
    \begin{block}{普通块 \texttt{block}}
        用于一般性陈述，标题与正文均可自定义。
    \end{block}
    \begin{alertblock}{警示块 \texttt{alertblock}}
        用于强调注意事项或需要引起警惕的内容。
    \end{alertblock}
    \begin{exampleblock}{示例块 \texttt{exampleblock}}
        用于举例、练习或代码说明。
    \end{exampleblock}
\end{frame}

\section{数学与定理}

\subsection{数学公式}

\begin{frame}{数学公式}
    行内公式如 \(e^{i\pi}+1=0\)；无编号行间公式用 \texttt{equation*}：
    \begin{equation*}
        J(\theta) = \mathbb{E}_{\pi_\theta}[G_t]
        = \sum_{s\in\mathcal{S}} d^\pi(s)\sum_{a\in\mathcal{A}}
        \pi_\theta(a|s)Q^\pi(s,a)
    \end{equation*}
    需要编号和交叉引用时用 \texttt{equation}：
    \begin{equation}
        \int_{-\infty}^{+\infty} e^{-x^2}\,\mathrm{d}x = \sqrt{\pi}
        \label{eq:gauss}
    \end{equation}
    式~\eqref{eq:gauss} 可通过 \texttt{\textbackslash eqref} 引用。
\end{frame}

\begin{frame}{多行公式}
    \begin{columns}[t]
        \begin{column}{0.48\textwidth}
            \texttt{align}：多行对齐，\texttt{\&} 为对齐点
            \begin{align}
                \nabla\cdot\mathbf{E}  & = \frac{\rho}{\varepsilon_0} \\
                \nabla\times\mathbf{B} & = \mu_0\mathbf{J}
                + \mu_0\varepsilon_0\frac{\partial\mathbf{E}}{\partial t}
            \end{align}
        \end{column}
        \begin{column}{0.48\textwidth}
            \texttt{multline}：单个公式跨行折行，首行靠左、末行靠右
            \begin{multline}
                a+b+c+d+e+f \\
                +g+h+i+j+k \\
                +l+m+n
            \end{multline}
        \end{column}
    \end{columns}
\end{frame}

\subsection{定理环境}

\begin{frame}{定理环境}
    Beamer 预定义了 \texttt{theorem}、\texttt{definition}、
    \texttt{lemma}、\texttt{proof} 等环境：
    \begin{theorem}[Cauchy--Schwarz 不等式]
        对任意向量 \(\mathbf{x},\mathbf{y}\in\mathbb{R}^n\)，有
        \((\sum_{i=1}^{n} x_i y_i)^2 \le (\sum_{i=1}^{n} x_i^2)(\sum_{i=1}^{n} y_i^2)\)，
        等号成立当且仅当 \(\mathbf{x}\) 与 \(\mathbf{y}\) 线性相关。
    \end{theorem}
    \begin{proof}
        对任意 \(t\in\mathbb{R}\)，\(\sum_i (x_i t + y_i)^2 \ge 0\) 展开为
        关于 \(t\) 的二次式，其判别式 \(\Delta = 4(\sum_i x_i y_i)^2 - 4\sum_i x_i^2 \sum_i y_i^2 \le 0\)，即得结论。
    \end{proof}
\end{frame}

\section{表格与插图}

\subsection{三线表}

\begin{frame}{三线表}
    表格推荐使用 \texttt{booktabs} 宏包的三线表风格：
    \begin{table}[h]
        \centering
        \caption{不同引擎的中文支持对比}
        \label{tab:engines}
        \begin{tabular}{lcc}
            \toprule
            引擎        & Unicode 支持 & 系统字体 \\
            \midrule
            pdf\LaTeX & 需 CJKutf8  & 否    \\
            Xe\LaTeX  & 原生         & 是    \\
            Lua\LaTeX & 原生         & 是    \\
            \bottomrule
        \end{tabular}
    \end{table}
    如表~\ref{tab:engines} 所示，本模板采用 Lua\LaTeX{}。
\end{frame}

\subsection{插图与分栏}

\begin{frame}{插图与分栏}
    \begin{columns}
        \begin{column}{0.45\textwidth}
            \begin{figure}[htpb]
                \centering
                \includegraphics[width=0.8\linewidth]{img/Example/SimpleDistributedStorageArchitecture.png}
                \caption{插图示例}
                \label{fig:example}
            \end{figure}
        \end{column}
        \begin{column}{0.5\textwidth}
            \begin{itemize}
                \item 分栏使用 \texttt{columns} 环境
                \item 矢量图（pdf、eps）优先于位图（png、jpg）
                \item \texttt{\textbackslash caption} 添加编号图题，
                      \texttt{\textbackslash ref} 引用，见图~\ref{fig:example}
            \end{itemize}
        \end{column}
    \end{columns}
\end{frame}

\subsection{TikZ 绘图}

\begin{frame}{TikZ 绘图}
    \begin{columns}
        \begin{column}{0.5\textwidth}
            \centering
            \begin{tikzpicture}[
                    neuron/.style={circle, draw=peaking, thick,
                            minimum size=6mm, inner sep=0},
                    label/.style={font=\footnotesize, align=center}
                ]
                % 各层神经元 / Neurons of each layer
                \foreach \i in {1,...,4}
                \node[neuron, fill=peaking!10] (I\i) at (0,{1.8-\i*1.0}) {};
                \foreach \i in {1,...,5}
                \node[neuron, fill=peaking!30] (H1\i) at (1.8,{2.0-\i*0.8}) {};
                \foreach \i in {1,...,5}
                \node[neuron, fill=peaking!30] (H2\i) at (3.6,{2.0-\i*0.8}) {};
                \foreach \i in {1,...,2}
                \node[neuron, fill=peaking!70, text=white] (O\i) at (5.4,{1.0-\i*1.4}) {};
                % 全连接边 / Fully connected edges
                \foreach \i in {1,...,4}
                \foreach \j in {1,...,5}
                \draw[peaking!20] (I\i) -- (H1\j);
                \foreach \i in {1,...,5}
                \foreach \j in {1,...,5}
                \draw[peaking!30] (H1\i) -- (H2\j);
                \foreach \i in {1,...,5}
                \foreach \j in {1,...,2}
                \draw[peaking!40] (H2\i) -- (O\j);
                % 层标签 / Layer labels
                \node[label] at (0,2.1)   {输入层\\\(\mathbf{x}\in\mathbb{R}^4\)};
                \node[label] at (1.8,2.1) {隐藏层 1\\\(\sigma(W_1\mathbf{x})\)};
                \node[label] at (3.6,2.1) {隐藏层 2\\\(\sigma(W_2\mathbf{h}_1)\)};
                \node[label] at (5.4,2.1) {输出层\\\(\hat{\mathbf{y}}\)};
            \end{tikzpicture}
        \end{column}
        \begin{column}{0.45\textwidth}
            \begin{itemize}
                \item 主题已加载 \texttt{tikz}，可直接绘图
                \item \texttt{\textbackslash foreach} 循环批量生成神经元与连接
                \item 颜色引用主题色 \texttt{peaking} 及其深浅变体
                \item 复杂图形建议单独维护 \texttt{.tikz} 文件
            \end{itemize}
        \end{column}
    \end{columns}
\end{frame}

\section{代码与参考文献}

\subsection{代码高亮}

\begin{frame}[fragile]{代码高亮}
    \texttt{listings} 已在主题中配置好样式，直接插入代码：
    \begin{columns}
        \begin{column}{0.48\textwidth}
            \begin{lstlisting}[language=C]
#include <stdio.h>
int main(void) {
    printf("Hello, SDU!\n");
    return 0;
}
\end{lstlisting}
        \end{column}
        \begin{column}{0.48\textwidth}
            \begin{lstlisting}[language=Python]
def fib(n):
    a, b = 0, 1
    for _ in range(n):
        a, b = b, a + b
    return a
\end{lstlisting}
        \end{column}
    \end{columns}
\end{frame}

\subsection{参考文献}

\begin{frame}{引用}
    \begin{itemize}
        \item 条目写在 \texttt{ref.bib} 中，用
              \texttt{\textbackslash cite\{key\}} 引用，
              例如 \cite{texbook} 与 \cite{beamerug}
        \item 软件与模板类资源可用 \texttt{@misc}/\texttt{@manual} 条目
              记录仓库地址，见 \cite{thu,pku,origin,sdulogo}
        \item 使用 \texttt{numeric-comp} 引用样式，由 \texttt{biber} 处理，
              引用标号以主题色显示
    \end{itemize}
\end{frame}

\begin{frame}[allowframebreaks]{参考文献}
    \printbibliography
\end{frame}

% 致谢页 / Thanks page
\begin{frame}
    \vfill
    {\Huge\bfseries Thanks!\par}
    \vskip 0.6em
        {\color{peaking}\rule{4em}{1.4pt}\par}
    \vfill
\end{frame}

\end{document}
